\documentclass[11pt]{article}

\usepackage[a4paper,margin=1in]{geometry}
\usepackage{amsmath,amssymb}
\usepackage{booktabs}
\usepackage{enumitem}

\title{Partial Assembly by Projection:\\
       an inexact matrix-free apply for the SFEM code generator}
\author{}
\date{}

\newcommand{\R}{\mathbb{R}}
\newcommand{\refcell}{\widehat{K}}
\newcommand{\Wbar}{\overline{W}}
\newcommand{\Sbar}{\overline{S}}
\newcommand{\xhat}{\boldsymbol{\xi}}

\begin{document}
\maketitle

\begin{abstract}
The matrix-free apply of a nonlinear elasticity operator carries a quadrature
loop because the material tangent varies across each element.  Projecting that
tangent onto the constants separates the loop into a small per-element object
and a tensor that belongs to the reference element alone.  The reference tensor
can then be integrated symbolically once, at generation time, and it turns out
to be far smaller than its size suggests: it is a Gram matrix, so its rank is
the dimension of the span of the element's gradient functions.  For a linear
simplex that rank is one, and the construction reduces exactly to the loperand.
This note states the technique, the compression, what the code generator emits,
how it vectorises, and the measurements taken so far.
\end{abstract}

\section{The exact apply, and why it has a loop}

Let $K$ be an element with reference cell $\refcell$ and affine or isoparametric
map $\boldsymbol{x}(\xhat)$, Jacobian $J$, and $\mathrm{d}V = \det J \,
\mathrm{d}\xhat$.  For a hyperelastic energy with first Piola stress $P(F)$, the
matrix-free action of the second form on an increment $h$ is
\begin{equation}
  (\mathcal{H} h)_{i,p}
  = \int_K S_{ijkl}\bigl(F(\boldsymbol{x})\bigr)\,
    \frac{\partial h_k}{\partial x_l}\,
    \frac{\partial \varphi_p}{\partial x_j}\;\mathrm{d}V ,
  \qquad
  S_{ijkl} = \frac{\partial P_{ij}}{\partial F_{kl}} ,
  \label{eq:exact}
\end{equation}
with $i$ the output component, $p$ the test function, and $\varphi_p$ the basis.

Pull the derivatives back to the reference cell with
$\partial/\partial x_j = \sum_n J^{-1}_{nj}\, \partial/\partial \xi_n$ and
collect the geometry into the tangent:
\begin{equation}
  \Sbar_{ikmn}
  = \sum_{j,l} S_{ijkl}\, J^{-1}_{nj}\, J^{-1}_{ml}\, \det J ,
  \label{eq:pullback}
\end{equation}
so that \eqref{eq:exact} becomes
\begin{equation}
  (\mathcal{H} h)_{i,p}
  = \int_{\refcell} \Sbar_{ikmn}(\xhat)\,
    \frac{\partial h_k}{\partial \xi_m}\,
    \frac{\partial \varphi_p}{\partial \xi_n}\;\mathrm{d}\xhat .
  \label{eq:reference}
\end{equation}
This $\Sbar$ is the \emph{loperand}: the flux differentiated once and already
contracted with the geometry, so that the element vector is nothing but
differences of its contractions.  It is the same object SFEM's hand-written
kernels build, and the same one \texttt{plans/loperand.py} isolates from a
material's flux.

Everything under the integral except $\Sbar$ is a property of the reference
element.  What forces a quadrature loop is that $\Sbar$ depends on $\xhat$,
because $F$ does.

\section{The projection}

Replace $\Sbar$ by its $L^2$ projection onto the constants on $K$, that is by
its element average $\Sbar^0 = |\refcell|^{-1}\int_{\refcell}\Sbar$.  It leaves
the integral and the sum separates:
\begin{equation}
  (\mathcal{H} h)_{i,p}
  \;\approx\;
  \Sbar^0_{ikmn}\, \sum_j h_{k,j}\, \Wbar_{jmpn} ,
  \qquad
  \Wbar_{jmpn}
  = \int_{\refcell}
    \frac{\partial \varphi_j}{\partial \xi_m}\,
    \frac{\partial \varphi_p}{\partial \xi_n}\;\mathrm{d}\xhat .
  \label{eq:projected}
\end{equation}
$\Sbar^0$ is one small object per element, computed once from the state.
$\Wbar$ belongs to the element and is integrated exactly, symbolically, at
generation time.  \emph{The quadrature loop leaves the kernel entirely.}

\subsection*{Where the approximation is, and where there is none}

On an affine simplex the deformation gradient is constant over the cell, so
$\Sbar$ is constant, so projecting it onto the constants returns it unchanged
and \eqref{eq:projected} is an identity rather than an approximation.  This
holds for a nonlinear material as much as a linear one: it is a statement about
the geometry, not about $P$.

That is the load-bearing fact of the whole construction, because it gives an
exact reference to test against.  On TRI3 and TET4 the projected kernel must
reproduce the exact one; only on curved or higher-order elements is there an
error, and there it must be measured rather than assumed.

\subsection*{Symmetry}

The element matrix entry implied by \eqref{eq:projected} is
$\Sbar^0_{ikmn}\Wbar_{jmpn}$.  Requiring it to be symmetric under exchanging the
two (component, node) pairs, and using $\Wbar_{jmpn} = \Wbar_{pnjm}$, forces
\begin{equation}
  \Sbar^0_{ikmn} = \Sbar^0_{kinm} .
  \label{eq:tangentsym}
\end{equation}
The pair that exchanges is $(i,n)$ against $(k,m)$: the output component with
the \emph{test} direction, and the increment component with the \emph{increment}
direction.  Packing on $(i,k)$ against $(m,n)$ instead is the plausible-looking
error, and it is invisible to any test that compares the construction against
itself.  Under \eqref{eq:tangentsym}, $\Sbar^0$ has $45$ independent components
in three dimensions and $10$ in two.

The derivation above assumed the element matrix is symmetric, and that is a
property of the \emph{formulation}, not of the method.  A material written as an
energy has a flux that is a gradient, so its tangent is a Hessian and
\eqref{eq:tangentsym} holds by equality of mixed partials.  A material written
as a residual has no such guarantee: its flux is not a gradient, its tangent is a
Jacobian, and it need not be symmetric at all.  Kelvin--Voigt viscosity is not
--- measured, $\max_{ijkl}|A_{ijkl}-A_{klij}| = 0.25$ --- so folding it into $45$
components averages its antisymmetric part away and yields a \emph{different
operator}: $5.3\times 10^{-2}$ relative error on TET4, an element where the
projection is provably exact.

The generator therefore decides symmetry per material, symbolically, and stores
$45$ or $81$ accordingly.  It defaults to unsymmetric, because the two errors are
not comparable: storing a full square for a symmetric tangent wastes memory,
while storing half of an unsymmetric one is silently wrong.

\section{Compression of the reference tensor}

\subsection{What it contains}

$\Wbar$ is integrated exactly, so its entries are rationals.  There are very few
of them, and on the simplices most are zero.

\begin{table}[ht]
\centering
\begin{tabular}{lrrrr}
\toprule
element & entries & distinct values & zeros & projection \\
\midrule
TRI3  &  36 &  3 &  20 & exact \\
TET4  & 144 &  3 & 108 & exact \\
QUAD4 &  64 &  6 &   0 & inexact \\
HEX8  & 576 & 10 &   0 & inexact \\
TET10 & 900 &  9 & 519 & inexact \\
\bottomrule
\end{tabular}
\caption{The symbolically integrated reference tensor.  At most ten distinct
rationals for any element.}
\label{tab:reference}
\end{table}

\subsection{Rank, which is where the structure really is}

The zeros are a symptom.  $\Wbar$, read as a matrix on the pair index $(j,m)$,
is the Gram matrix of the element's gradient functions
$\{\partial\varphi_j/\partial\xi_m\}$ in $L^2(\refcell)$.  Its rank is therefore
the dimension of the span of those functions and nothing more.

\begin{table}[ht]
\centering
\begin{tabular}{lrrl}
\toprule
element & order & rank & span of the gradient functions \\
\midrule
TRI3  &  6 & 1 & constants \\
TET4  & 12 & 1 & constants \\
QUAD4 &  8 & 3 & $1,\ \xi_0,\ \xi_1$ \\
HEX8  & 24 & 7 & $1$ and the six bilinear monomials \\
TET10 & 30 & 4 & $1,\ \xi_0,\ \xi_1,\ \xi_2$ \\
\bottomrule
\end{tabular}
\caption{Rank of the reference tensor, which equals the dimension of the span
of the element's gradient functions exactly.}
\label{tab:rank}
\end{table}

A linear simplex is rank one because its gradients are constants.  Rank one is
precisely the loperand of Section~1: one object per element, contracted once.
So the rank factorisation is not a second idea beside the loperand; it is the
same structure carried to elements whose gradients are not constant.

Factor exactly over the rationals,
\begin{equation}
  \Wbar = C^{\mathsf T} X C ,
  \qquad C \in \mathbb{Q}^{r \times N d},
  \qquad X \in \mathbb{Q}^{r \times r},
\end{equation}
with $r$ the rank, $N$ the node count and $d$ the dimension.  The factorisation
is checked against the tensor it factors when it is built.

\section{The generated kernel}

\subsection{Staging}

Substituting the factorisation into \eqref{eq:projected} and naming each
intermediate gives four stages:
\begin{align}
  P_{kam} &= \sum_j C_{a,(j,m)}\, h_{k,j}
            && \text{the increment, compressed} \\
  Y_{ian} &= \sum_{k,m} \Sbar^0_{ikmn}\, P_{kam}
            && \text{the tangent applied} \\
  Q_{ibn} &= \sum_a X_{ab}\, Y_{ian}
            && \text{the middle factor} \\
  \text{out}_{ip} &= \sum_{b,n} C_{b,(p,n)}\, Q_{ibn}
            && \text{expanded back to nodes}
\end{align}
Every sum skips its zero coefficients, so a zero of $C$ or $X$ never enters an
expression rather than being cancelled out of one later.

The staging is the point and not an implementation detail.  Expanding the four
stages into one expression per output and letting common-subexpression
elimination rediscover the structure returns the dense operation count exactly;
the saving lives in keeping the intermediates named.

\subsection{Choosing the contraction}

Section~5 stages the action through the rank factorisation.  That is not always
the cheaper arrangement.  The factorisation has to compress the increment and
expand the result, and where $\Wbar$'s rank is a large fraction of its order,
those stages cost more than they save.  The alternative contracts $\Wbar$ against
the increment directly,
\begin{align}
  G_{kmpn} &= \sum_j \Wbar_{jmpn}\, h_{k,j} \\
  \text{out}_{ip} &= \sum_{k,m,n} \Sbar^0_{ikmn}\, G_{kmpn} ,
\end{align}
in which $\Sbar^0$ appears once, in a single $27$-term contraction per output,
and there is nothing to compress or expand.

Which wins is a property of the element, so the generator measures both and
emits the cheaper.  Counted after common-subexpression elimination:

\begin{table}[ht]
\centering
\begin{tabular}{lrrrl}
\toprule
element & staged & gradient-first & rank & chosen \\
\midrule
TET4  &  207 &  342 &  1 of 12 & staged \\
TET10 & 1380 & 1605 &  4 of 30 & staged \\
HEX8  & 2610 & 1731 &  7 of 24 & gradient-first \\
\bottomrule
\end{tabular}
\caption{Operations for the two contraction orderings.  HEX8's rank is high
enough that the factorisation stops paying.}
\end{table}

On HEX8 this takes the emitted apply from $2732$ operations to $1801$, against a
hand-written \texttt{hex8\_SdotHdotG} at $1558$, and lifts its measured rate from
about $100$ to about $132$~MDOF/s.

The counting has to be done \emph{after} elimination.  Measured on the
unsimplified expressions, gradient-first looks worse on every element, which is
how it was first dismissed.

\subsection{Fold the constants, do not tabulate them}

The entries of $\Wbar$ are known at generation time, so they are emitted as
literals.  The tempting alternative --- store the distinct values in a runtime
table and index it --- compresses the memory and destroys the arithmetic: the
constants stop being constants, the compiler can no longer fold them, and every
multiply by a zero entry survives into the emitted code.

This is not hypothetical.  The existing prototype
(\texttt{python/codegen/neohookean\_partial\_assembly.py}) emits, for TET4, a
table holding exactly $1$, $-1$ and $0$, and its kernels then multiply by the
zero entry twelve times each; the reported ``$97.9\%$ memory saving'' buys
$141$ stored numbers that were $0$ and $\pm 1$ at a cost of several hundred
floating-point operations.

\begin{table}[ht]
\centering
\begin{tabular}{lrrrr}
\toprule
element & rank/order & dense, folded & staged, folded & runtime table \\
\midrule
TRI3  & 1/6  &   96 &  \textbf{52} &  244 \\
TET4  & 1/12 &  486 & \textbf{207} & 1467 \\
QUAD4 & 3/8  &  238 &          238 &  431 \\
HEX8  & 7/24 & 4059 & \textbf{2610} & 6330 \\
TET10 & 4/30 & 4018 & \textbf{1380} & 8349 \\
\bottomrule
\end{tabular}
\caption{Operations in the element apply, counted after
common-subexpression elimination.  Preliminary: these are symbolic operation
counts, not measured run times.}
\label{tab:cost}
\end{table}

Table~\ref{tab:cost} says where the staging pays.  The saving tracks the rank
against the order: TET10 at $4/30$ gains $2.9\times$ and TET4 at $1/12$ gains
$2.4\times$, while QUAD4 at $3/8$ comes out exactly even --- not enough
structure to beat common-subexpression elimination on the dense form.  Against
the runtime-table form the improvement is $7.1\times$ on TET4 and $6.1\times$ on
TET10.

\section{Vectorisation}

The generated kernels process a block of \texttt{VECTOR\_SIZE} elements at a
time, with the lane loop innermost and marked \verb|#pragma omp simd|.  The
staged form suits that layout well, and for reasons worth stating separately.

\begin{description}[leftmargin=1.5em]
\item[The reference tensor is lane-invariant.]  $C$ and $X$ are compile-time
  rationals, identical for every element and therefore for every lane.  Folded,
  they appear as immediates: no load, no broadcast, and a multiply by $\pm 1$ or
  by a power of two is folded away entirely.  Behind a runtime table the same
  values become a memory reference whose address does not depend on the lane ---
  a broadcast load rather than an immediate --- and, worse, an opaque value that
  no longer permits folding.  This is the arithmetic half of the cost in
  Table~\ref{tab:cost}.

\item[The tangent is a per-element stream.]  $\Sbar^0$ is $45$ numbers per
  element in three dimensions, stored one component-array per index in the
  structure-of-arrays layout the framework already uses for geometry.  Each
  component is then a unit-stride read across the lane loop, which vectorises
  without gather.

\item[The intermediates are lane-width vectors.]  $P$, $Y$ and $Q$ are
  per-element scalars, so in a blocked kernel each becomes one vector register.
  Their counts are the stage sizes in Table~\ref{tab:stages}, and they are the
  binding constraint: on HEX8, $63 + 63 + 63$ live intermediates at a lane width
  of eight doubles will not stay in registers, and the kernel will spill.  Rank
  is what controls this, since all three stage sizes are $d^2 r$.

\item[The output scatter is unchanged.]  The last stage writes one value per
  node per component, exactly as the existing apply kernels do, so it reuses the
  same scatter --- atomic on the plain traversal, and plain accumulation into
  thread-private scratch on the packed one.
\end{description}

\begin{table}[ht]
\centering
\begin{tabular}{lrrrrr}
\toprule
element & rank & $|P|$ & $|Y|$ & $|Q|$ & outputs \\
\midrule
TRI3  & 1 &  4 &  4 &  4 &  6 \\
TET4  & 1 &  9 &  9 &  9 & 12 \\
QUAD4 & 3 & 12 & 12 & 12 &  8 \\
HEX8  & 7 & 63 & 63 & 63 & 24 \\
TET10 & 4 & 36 & 36 & 36 & 30 \\
\bottomrule
\end{tabular}
\caption{Live intermediates per element in each stage.  In a blocked kernel each
is one vector register, so these are the register-pressure figures.}
\label{tab:stages}
\end{table}

A consequence worth drawing out: on the elements where the projection is exact
the register pressure is also lowest, and on HEX8 --- where the projection is an
approximation --- it is highest.  The two difficulties arrive together, and they
are the same difficulty: rank controls both.  This is why HEX8 is also the
element on which the staged form loses to the direct contraction, which needs no
intermediates at all.  A rank-truncated factorisation (keeping $r' < r$ modes)
remains the knob for elements where staging still wins, at the price of a second
and larger approximation.

\section{Verification}

Each stage can be wrong in a way the next cannot see, so the construction is
checked at four independent points.

\begin{enumerate}[leftmargin=1.5em]
\item \textbf{The symbolic integration}, against four-point Gauss--Legendre
  quadrature on QUAD4 and HEX8, entry by entry: agreement to $10^{-16}$.
\item \textbf{The factorisation}, against the tensor it factors:
  $C^{\mathsf T} X C = \Wbar$ exactly, over the rationals.
\item \textbf{The action}, against \eqref{eq:exact} integrated directly on a
  concrete tetrahedron, for linear elasticity and for neohookean Ogden.
\item \textbf{The generated operator}, inside SFEM, driven as a solver drives it:
  \texttt{initialize}, \texttt{apply}, \texttt{inexact\_update},
  \texttt{inexact\_apply}.
\end{enumerate}

Only the third can catch a wrong index convention, and it did: the symmetry
\eqref{eq:tangentsym} had been implemented as $(i,k)$ against $(m,n)$, which the
first two checks and any comparison of the staged form against the dense form
all pass, because they share the error.

Only the fourth can catch anything about integration.  It found that the emitted
operator source omitted the C~ABI header, leaving every kernel signature
unparseable; that the affine geometry cache the split reads was not built unless
some \emph{other} path had asked for it; and that \texttt{inexact\_supported()}
reported support on meshes where the geometry cache cannot exist.  None of these
is visible in the emitted text.

\subsection*{How large is the projection error where it is not zero}

Section~3 leaves the curved case open: on TET10 and HEX8 the projection is an
approximation and its error has to be measured.  It is measured by warping the
displacement --- an increasing perturbation of the reference configuration,
checked at each step to keep $\det F > 0$ --- and comparing the projected
neohookean apply against the exact one at each severity.

\begin{table}[ht]
\centering
\begin{tabular}{lrrrr}
\toprule
warp $w$ & 0 & 0.4 & 1.0 & 1.6 \\
\midrule
TET4  & 1e-15 & 1e-15 & 1e-15 & 1e-15 \\
TET10 & 2.4e-14 & 1.0e-02 & --- & 4.3e-02 \\
HEX8  & 2.7e-15 & --- & --- & 3.7e-04 \\
\bottomrule
\end{tabular}
\caption{Relative error of the projected apply against the exact one, $n=16$.
The $w=0$ column is the control: it recovers the exactness proof of Section~3
and shows the sweep is measuring curvature and not the harness.  TET4 is flat
across the sweep, as it must be.}
\end{table}

TET10 is two orders worse than HEX8 at comparable warp.  Both are quadratic in
the geometry, but TET10 carries a quadratic \emph{displacement} whose gradient
varies linearly over the cell, and it is that variation the projection onto the
constants discards.  The error is a real property of the approximation, not a
defect: it falls under refinement at $O(h^{1.07})$ (see Methodology for why the
first measurement of this appeared not to converge), so the projected operator
converges to the exact one at the rate the discretisation itself does.

\section{Methodology}

Several of the numbers in earlier drafts of this note were wrong.  The errors
were not in the kernels, and they are worth recording, because each is a way of
measuring the wrong thing while appearing to measure the right one.

\paragraph{Time the kernel, not the harness.}  The first throughput figures put a
serial \texttt{std::fill} of the output arrays inside the timed region, and
wrapped a single call per sample so that each paid OpenMP team startup and a cold
cache.  Neither is part of the operator.  The distortion is invisible at one
thread and grows with the thread count --- the worst possible shape for a scaling
study: the single-thread figures barely moved when it was fixed ($2.52$ to
$2.68$~MDOF/s) while the eight-thread figures moved by half ($13.04$ to
$18.75$).  Reported scaling was $5.4\times$ where it is really $7.0\times$.

\paragraph{A ratio can be flat while the thing it measures converges.}
Measured against a \emph{smooth} increment, the TET10 projection error sat at
$1.0\times10^{-2}$ across an eightfold refinement and looked like a defect.  It
is not.  For a smooth field the assembled $(Hh)_p$ is a discrete second
derivative, so neighbouring elements cancel and the denominator collapses by an
extra order, while the per-element projection errors carry independent signs and
do not cancel at all.  Against a white-noise increment the same kernel converges
at $O(h^{1.07})$.  Before concluding that an approximation fails to converge,
re-measure it with a random vector.

\paragraph{Do not measure at a convenient mesh size.}  The stored tangent is
\texttt{metric\_tensor\_t}, which is \texttt{float}.  On a mesh whose spacing is
a power of two the tangent entries are exactly representable and the difference
from the exact apply collapses to $10^{-14}$, which looks like proof of something
far stronger than is true.  The same kernels at resolution $12$:

\begin{center}
\begin{tabular}{lrrr}
\toprule
element & $n=8$ & $n=12$ & $n=16$ \\
\midrule
TET4  & 6.1e-16 & 1.2e-07 & 1.8e-15 \\
TET10 & 1.4e-14 & 1.0e-05 & 4.8e-14 \\
HEX8  & --      & 6.8e-07 & --      \\
\bottomrule
\end{tabular}
\end{center}

So the split reproduces the exact apply to the precision of its store, and that
precision is element-dependent: TET10's quadratic basis gives its tangent a wider
dynamic range and costs about two more digits than TET4 for the same
\texttt{float} store.  That is an input to choosing a store precision, not a
tolerance to widen.

\paragraph{A conclusion from one machine is a conclusion about that machine.}
On an Apple laptop every column fell by a fifth past eight threads, and the
atomic scatter looked like the scalability limit.  It is not: the laptop has
eight performance cores and two efficiency cores, \texttt{schedule(static)} hands
them equal chunks, and the loop waits for the slow ones.  On $72$ homogeneous
Grace cores the same kernels scale essentially linearly.

\section{Measured performance}

Two frames, and they are not interchangeable.  The first calls the generated
kernels directly with the geometry precomputed; the second goes through
\texttt{Op::apply}, which is the path a solve actually takes.  At $170{,}373$
degrees of freedom the same kernel reads about $5$~MDOF/s in the first and
$1.46$ in the second.  Only the second is comparable with the library.

\subsection*{Against the library's own operators}

Measured with SFEM's \texttt{bench\_op}, a spike copy of which reports setup and
application separately.  \texttt{build64}, eight threads, ${\sim}10.4$M~dof.

\begin{table}[ht]
\centering
\begin{tabular}{lrrrr}
\toprule
element & hand-written PA & exact generated & generated split & vs.\ hand-written \\
\midrule
TET10 & 137.4 & 27.2 & 219.6 & $1.60\times$ faster \\
HEX8  & 162.9 & 47.5 & 100.2 & $1.63\times$ behind \\
TET4  & (none) & 47.9 & 116.9 & --- \\
\bottomrule
\end{tabular}
\caption{MDOF/s for the apply.  ``Hand-written PA'' is
\texttt{NeoHookeanOgden}, which enables partial assembly by hand for HEX8 and
TET10 and not for TET4 --- which is exactly where the generated kernels were
ahead and behind before this work.}
\end{table}

The assemblies, measured separately: TET10 $0.392$\,s for the generated split
against $0.653$\,s hand-written; HEX8 level at $0.414$ against $0.427$.  So on
TET10 the generated split beats hand-written partial assembly on both halves.
On HEX8 it doubles the exact generated kernel but remains behind the
hand-written one, and since the assemblies are level the whole gap is the
contraction.

Break-even needs a reference that performs no setup of its own.  The
hand-written operator assembles too, so the exact generated apply is the honest
denominator: against it, TET10 breaks even at $1.03$ applies per tangent.

\subsection*{Threads}

\begin{table}[ht]
\centering
\begin{tabular}{lrrrrr}
\toprule
 & 1 & 8 & 32 & 64 & 72 \\
\midrule
Apple, exact      &  2.68 & 18.75 & --- & --- & --- \\
Apple, stored f16 &  7.61 & 55.15 & --- & --- & --- \\
Grace, exact      &  1.33 & 10.68 & 41.65 &  82.13 &  92.28 \\
Grace, stored f16 &  2.85 & 25.30 & 97.67 & 191.62 & 209.60 \\
\bottomrule
\end{tabular}
\caption{MDOF/s at $206{,}763$~dof, Mooney--Rivlin with Kelvin--Voigt viscosity.
Grace scales $69\times$ on $72$ cores for the exact apply and $73\times$ for the
stored one.  The Apple figures peak at eight threads, not ten.}
\end{table}

On Grace \texttt{fp16} is not worth carrying: at one thread it is slower than
both \texttt{fp32} and \texttt{fp64}, and at $72$ threads it is within two per
cent of \texttt{fp32} while costing four decimal digits.  \texttt{fp32} is the
default on both machines.

\section{Status}

The split is implemented and reachable from SFEM.  \texttt{Op} carries
\texttt{inexact\_supported()}, \texttt{inexact\_update()} and
\texttt{inexact\_apply()}, with defaults that refuse, so no existing operator or
caller changes behaviour.  \texttt{inexact\_update} is explicit by design: the
stored tangent stays valid until the next call, which is the reuse the split
exists for, and \texttt{inexact\_apply} takes no state at all, so it cannot
silently use a stale one.

Three runtime-typed C~ABI entry points per material --- assemble, apply from a
\texttt{metric\_tensor\_t} store, apply from a compressed one --- dispatch on
\texttt{smesh::ElemType} with the scalar type as a \texttt{smesh::PrimitiveType}.
\texttt{linear\_elasticity} is generated in-tree with the split enabled, and its
operator reproduces its own exact apply to the store's precision on TET4, TET10
and HEX8.

Open:
\begin{enumerate}[leftmargin=1.5em]
\item The 2D kernels are generated but unreachable: smesh's adjugate fill has no
  TRI3 or QUAD4 case, so the affine geometry cache the split reads cannot be
  built.  The operator reports \texttt{inexact\_supported() == false} there
  rather than failing at run time, and will start working by itself if smesh
  gains 2D support.
\item HEX8's contraction is still behind the hand-written kernel by about
  $1.6\times$, entirely in the contraction rather than the assembly.
\item The 2D store size is emitted and has never executed.
\item A permanent test belongs in \texttt{frontend/tests/}, beside
  \texttt{sfem\_PartialAssemblyTest.cpp}, rather than in a spike.
\end{enumerate}

\end{document}
